% GENERATED FILE. Edit canonical lessons, then run npm run generate:books.
% canonical-source-hash: fnv1a64:10f63f861f9bae61
% canonical-lessons: RU-C439-povorot, RU-C439-razvilka, RU-C439-tupik, RU-C439-srezat, RU-C439-ogibat

\chapter{Turns and dead ends}
\label{ch:ru-turns-and-dead-ends}
\hlchaptermodality{439}

\begin{chapteropening}
\textbf{By the end of this chapter:} \emph{I can follow a route: a turn, a fork in the road, a dead end, taking a short cut and going round something.}

Complete the last lesson of chapter 439: I can follow a route: a turn, a fork in the road, a dead end, taking a short cut and going round something.
\end{chapteropening}

\section[\texorpdfstring{\ru{поворот} (povorót)}{поворот (povorót)}]{\ru{поворот} (povorót) — a turn, a bend}

\label{lesson:RU-C439-povorot}



\begin{quote}
Before the new one: say the Russian for a landmark, then the Russian for a route map.
\end{quote}

\subsection*{You'll want to know: \ru{поворот}}
\textbf{\ru{поворот}} (\emph{povorót}) — \textquotedblleft{}a turn, a bend\textquotedblright{}.

\begin{quote}
\ru{После} \ru{поворота} \ru{будет} \ru{мост}.
\end{quote}

\begin{quote}
\emph{Pósle povoróta búdet most.}
\end{quote}

\begin{quote}
After the bend there will be a bridge.
\end{quote}

Masculine. \emph{\ru{Поворот} \ru{направо}} — a right turn.

\begin{cousinweb}
\emph{\ru{поворот}} is built from \emph{\ru{по}-} and \emph{\ru{ворот}} 'a turning', which EtymDB takes from Proto-Slavic \emph{vortŭ}, on the root of \emph{\ru{вертеть}} 'to turn'.
\end{cousinweb}

\begin{grammarlens}[title={What you've built}]
One more word for telling, explaining and describing at B1.
\end{grammarlens}

\subsection*{Your turn}
\begin{itemize}
\raggedright
  \item \emph{Say it:} \emph{povorót}
  \item \emph{Say it:} \emph{povorót}, once more
  \item \emph{Say it:} a sentence of your own with \emph{povorót}
  \item \emph{Recall:} say the Russian for a landmark, then the Russian for a route map, then say \emph{povorót} again
\end{itemize}

\begin{tcolorbox}[breakable,colback=teal!4,colframe=teal!35!black,title={Before you move on}]
What does \ru{поворот} mean? (\textquotedblleft{}A turn, a bend\textquotedblright{}.) Say it once more.
\end{tcolorbox}
\section[\texorpdfstring{\ru{развилка} (razvílka)}{развилка (razvílka)}]{\ru{развилка} (razvílka) — a fork (in a road)}

\label{lesson:RU-C439-razvilka}



\begin{quote}
Before the new one: say the Russian for a route map, then the Russian for a turn.
\end{quote}

\subsection*{You'll want to know: \ru{развилка}}
\textbf{\ru{развилка}} (\emph{razvílka}) — \textquotedblleft{}a fork (in a road)\textquotedblright{}.

\begin{quote}
\ru{На} \ru{развилке} \ru{мы} \ru{повернули} \ru{налево}.
\end{quote}

\begin{quote}
\emph{Na razvílke my povernúli nalévo.}
\end{quote}

\begin{quote}
At the fork we turned left.
\end{quote}

Feminine. Built on \emph{\ru{вилка}} — a fork.

\begin{cousinweb}
\emph{\ru{развилка}} is built from \emph{\ru{раз}-} 'apart' and \emph{\ru{вилка}} 'a fork', from \emph{\ru{вилы}} 'a pitchfork'. EtymDB lists \emph{\ru{вилка}} without a deeper source.
\end{cousinweb}

\begin{grammarlens}[title={What you've built}]
One more word for telling, explaining and describing at B1.
\end{grammarlens}

\subsection*{Your turn}
\begin{itemize}
\raggedright
  \item \emph{Say it:} \emph{razvílka}
  \item \emph{Say it:} \emph{razvílka}, once more
  \item \emph{Say it:} a sentence of your own with \emph{razvílka}
  \item \emph{Recall:} say the Russian for a route map, then the Russian for a turn, then say \emph{razvílka} again
\end{itemize}

\begin{tcolorbox}[breakable,colback=teal!4,colframe=teal!35!black,title={Before you move on}]
What does \ru{развилка} mean? (\textquotedblleft{}A fork (in a road)\textquotedblright{}.) Say it once more.
\end{tcolorbox}
\section[\texorpdfstring{\ru{тупик} (tupík)}{тупик (tupík)}]{\ru{тупик} (tupík) — a dead end, a cul-de-sac; an impasse}

\label{lesson:RU-C439-tupik}



\begin{quote}
Before the new one: say the Russian for a turn, then the Russian for a fork.
\end{quote}

\subsection*{You'll want to know: \ru{тупик}}
\textbf{\ru{тупик}} (\emph{tupík}) — \textquotedblleft{}a dead end, a cul-de-sac; an impasse\textquotedblright{}.

\begin{quote}
\ru{Эта} \ru{улица} — \ru{тупик}.
\end{quote}

\begin{quote}
\emph{Éta úlitsa — tupík.}
\end{quote}

\begin{quote}
This street is a dead end.
\end{quote}

Masculine; the stress moves: \emph{\ru{тупика}} (\emph{tupiká}). Figurative: \emph{\ru{в} \ru{тупике}} — at a dead end.

\begin{cousinweb}
\emph{\ru{тупик}} is built from \emph{\ru{тупой}} 'blunt' and the ending \emph{-\ru{ик}} (EtymDB): a street with a blunt end.
\end{cousinweb}

\begin{grammarlens}[title={What you've built}]
One more word for telling, explaining and describing at B1.
\end{grammarlens}

\subsection*{Your turn}
\begin{itemize}
\raggedright
  \item \emph{Say it:} \emph{tupík}
  \item \emph{Say it:} \emph{tupík}, once more
  \item \emph{Say it:} a sentence of your own with \emph{tupík}
  \item \emph{Recall:} say the Russian for a turn, then the Russian for a fork, then say \emph{tupík} again
\end{itemize}

\begin{tcolorbox}[breakable,colback=teal!4,colframe=teal!35!black,title={Before you move on}]
What does \ru{тупик} mean? (\textquotedblleft{}A dead end, a cul-de-sac; an impasse\textquotedblright{}.) Say it once more.
\end{tcolorbox}
\section[\texorpdfstring{\ru{срезать} (srézat)}{срезать (srézat)}]{\ru{срезать} (srézat) — to take a short cut; to cut off}

\label{lesson:RU-C439-srezat}



\begin{quote}
Before the new one: say the Russian for a fork, then the Russian for a dead end.
\end{quote}

\subsection*{You'll want to know: \ru{срезать}}
\textbf{\ru{срезать}} (\emph{srézat}) — \textquotedblleft{}to take a short cut; to cut off\textquotedblright{}.

\begin{quote}
\ru{Мы} \ru{срезали} \ru{дорогу} \ru{через} \ru{парк}.
\end{quote}

\begin{quote}
\emph{My srézali dorógu chérez park.}
\end{quote}

\begin{quote}
We took a short cut through the park.
\end{quote}

Perfective; the imperfective \emph{\ru{срезать}} (\emph{srezát}) has the stress on the \ru{а}. \emph{\ru{Срезать} \ru{путь}} — to take a short cut. Future \emph{\ru{срежу}, \ru{срежешь}}.

\begin{cousinweb}
\emph{\ru{срезать}} is built from \emph{\ru{с}-} 'off' and \emph{\ru{резать}} 'to cut', which is inherited from Proto-Slavic \emph{rězati} (EtymDB).
\end{cousinweb}

\begin{grammarlens}[title={What you've built}]
One more word for telling, explaining and describing at B1.
\end{grammarlens}

\subsection*{Your turn}
\begin{itemize}
\raggedright
  \item \emph{Say it:} \emph{srézat}
  \item \emph{Say it:} \emph{srézat}, once more
  \item \emph{Say it:} a sentence of your own with \emph{srézat}
  \item \emph{Recall:} say the Russian for a fork, then the Russian for a dead end, then say \emph{srézat} again
\end{itemize}

\begin{tcolorbox}[breakable,colback=teal!4,colframe=teal!35!black,title={Before you move on}]
What does \ru{срезать} mean? (\textquotedblleft{}To take a short cut; to cut off\textquotedblright{}.) Say it once more.
\end{tcolorbox}
\section[\texorpdfstring{\ru{огибать} (ogibát)}{огибать (ogibát)}]{\ru{огибать} (ogibát) — to go round, to skirt}

\label{lesson:RU-C439-ogibat}



\begin{quote}
Before the new one: say the Russian for a dead end, then the Russian for to take a short cut.
\end{quote}

\subsection*{You'll want to know: \ru{огибать}}
\textbf{\ru{огибать}} (\emph{ogibát}) — \textquotedblleft{}to go round, to skirt\textquotedblright{}.

\begin{quote}
\ru{Дорога} \ru{огибает} \ru{озеро}.
\end{quote}

\begin{quote}
\emph{Doróga ogibáyet ózero.}
\end{quote}

\begin{quote}
The road goes round the lake.
\end{quote}

Imperfective; perfective \emph{\ru{обогнуть}} (\emph{\ru{обогну}, \ru{обогнёшь}}).

\begin{cousinweb}
\emph{\ru{огибать}} is built from \emph{\ru{о}-} 'round' and \emph{-\ru{гибать}}, on the root of \emph{\ru{гнуть}} 'to bend', which is inherited from Proto-Slavic \emph{gŭnǫti} (EtymDB).
\end{cousinweb}

\begin{grammarlens}[title={What you've built}]
One more word for telling, explaining and describing at B1.
\end{grammarlens}

\subsection*{Your turn}
\begin{itemize}
\raggedright
  \item \emph{Say it:} \emph{ogibát}
  \item \emph{Say it:} \emph{ogibát}, once more
  \item \emph{Say it:} a sentence of your own with \emph{ogibát}
  \item \emph{Recall:} say the Russian for a dead end, then the Russian for to take a short cut, then say \emph{ogibát} again
\end{itemize}

\begin{tcolorbox}[breakable,colback=teal!4,colframe=teal!35!black,title={Before you move on}]
What does \ru{огибать} mean? (\textquotedblleft{}To go round, to skirt\textquotedblright{}.) Say it once more.
\end{tcolorbox}
